\chapter{The Problem: Multilingual Speech Recognition with Language Identification}
\chaptermark{The Problem}

\section{Two tasks, one model}

\begin{definition}{Automatic speech recognition (ASR)}{asr}
Let $\mathcal X$ be the set of finite audio signals and $\Sigma^*$ the set of finite strings over a character set $\Sigma$
(here: Latin letters plus the Devanagari, Bengali--Assamese, Telugu and Odia scripts). ASR is the task of learning a map
$h:\mathcal X\to\Sigma^*$ such that $h(\vx)$ is the words spoken in $\vx$, written in the conventional script of the language.
\end{definition}

\begin{definition}{Spoken language identification (LID)}{lid}
Given a finite set of languages $\mathcal L$, LID is the task of learning $g:\mathcal X\to\mathcal L$, or better a probability
vector $g(\vx)\in\Delta^{|\mathcal L|-1}$ (the probability simplex), whose largest entry is the language spoken in $\vx$.
\end{definition}

In this project $\mathcal L$ is \emph{closed}: the model always answers with one of the seven languages, even for French audio. A
closed-set classifier never says ``none of these''.

\paragraph{Why not two systems?} The classical pipeline is \emph{LID $\to$ route to a language-specific recogniser}. It has two
weaknesses:
\begin{enumerate}
\item \textbf{Error propagation.} If LID is wrong, the wrong recogniser runs and the transcript is garbage. The cost of this is
measured by the \emph{oracle-vs-predicted gap}: the CER when the true language is given minus the CER when the predicted language
is used (Section~\ref{sec:oracle}).
\item \textbf{No sharing.} Seven separate models learn seven times what speech sounds like. Hindi and Marathi share the Devanagari
script, Bengali and Assamese share a script, and all six Indic languages share many sounds and loan words.
\end{enumerate}
A single model with shared representations can transfer knowledge between related languages, and its LID decision is
taken by the same network that will transcribe, so the two cannot disagree about what was heard.

\section{The seven languages}

\begin{center}
\small
\begin{tabular}{lllll}
\toprule
Code & Language & Family & Script & ``computer'' as usually written\\
\midrule
hi & Hindi & Indo-Aryan & Devanagari & \deva{कंप्यूटर}\\
mr & Marathi & Indo-Aryan & Devanagari & \deva{संगणक} / \deva{कॉम्प्युटर}\\
bn & Bengali & Indo-Aryan (Eastern) & Bengali & \beng{কম্পিউটার}\\
as & Assamese & Indo-Aryan (Eastern) & Bengali--Assamese & \beng{কম্পিউটাৰ}\\
or & Odia & Indo-Aryan (Eastern) & Odia & \oriy{କମ୍ପ୍ୟୁଟର}\\
te & Telugu & Dravidian & Telugu & \telu{కంప్యూటర్}\\
en & English & Germanic & Latin & computer\\
\bottomrule
\end{tabular}
\end{center}

Three facts in this table drive many later decisions:
\begin{itemize}
\item Assamese and Bengali share a script (they differ in a few letters, e.g.\ Assamese \beng{ৰ} vs Bengali \beng{র}), and Hindi and
Marathi share Devanagari. LID between these pairs cannot rely on the script of the output; it must come from the audio.
\item Odia is closely related to Bengali and Assamese (Eastern Indo-Aryan). This is why the new Odia token is initialised from the
Bengali one (Chapter~\ref{ch:odia}).
\item Every Indic script is far from ASCII in Unicode. Whisper's tokenizer was built mostly on English text, so Indic text costs
many more tokens per second of speech (Chapter~\ref{ch:whisper}).
\end{itemize}

\paragraph{Code-switching.} In \emph{code-switched} speech the speaker mixes languages inside one sentence: ``\deva{मेरा} computer
\deva{बंद हो गया}''. We include Hindi--English and Bengali--English code-switched speech (the MUCS 2021 corpus) for training and
testing. A code-switched utterance is labelled with its \emph{matrix} language (hi or bn).

\section{Goals and scope decisions}

The goals came from the course proposal and the literature review (\file{ie643.md}):
\begin{enumerate}
\item One Transformer model that transcribes all seven languages and reports the language with a confidence.
\item Keep LID good after fine-tuning (Whisper's LID is known to collapse when it is fine-tuned on new languages).
\item Handle the data imbalance between languages.
\item Be robust to noise and do not invent text on silence or music (``hallucination'').
\item Compare honestly against the strongest open baselines: IndicConformer-600M (AI4Bharat), Qwen3-ASR (0.6B, 1.7B) and an
external LID model (SpeechBrain ECAPA, VoxLingua107).
\end{enumerate}

Scope decisions made during the project (they override the original slides):
\begin{itemize}
\item Final language set: hi, en, te, as, or, bn, mr. Tamil and a low-resource tier (Maithili, Konkani, Santali) were dropped.
\item About 90 hours per language instead of the full corpora (disk and compute budget).
\item English: LibriSpeech \texttt{train-clean-100} was added so that English had comparable hours.
\end{itemize}

\section{The four design principles from the literature}

\begin{intuition}[: what the literature review told us to do]
\begin{enumerate}
\item \textbf{Keep the language token as a training target.} If the decoder is always \emph{told} the language during fine-tuning,
it never practises \emph{predicting} it, and LID decays. Making the language token a supervised output keeps it trained.
\item \textbf{Add an auxiliary encoder LID loss} (LID-CTC). In ML-SUPERB 2.0 (Interspeech 2025) this gave a 14\% relative LID
improvement.
\item \textbf{Balance the data deterministically} (two-stage upsampling, as in Polyglot-Lion) rather than letting the biggest corpus
dominate.
\item \textbf{No denoiser in front of the recogniser.} Several 2025--26 studies found that speech enhancement in front of a strong
ASR model \emph{hurts}. Instead: training-time augmentation, and non-speech clips with empty transcripts (Canary, Calm-Whisper).
\end{enumerate}
\end{intuition}

\section{The hardware}
Everything ran on one NVIDIA Thor (aarch64, about 112--122 GB of \emph{unified} memory shared by CPU and GPU, CUDA 13). Unified
memory has one practical consequence that appears in the log: two Whisper trainings at once exhaust it and both stall, so all
trainings were run one after another.

\begin{vivatip}
If asked ``what is new here?'', the honest answer is: no single component is new; the contribution is the combination (Whisper
turbo + LoRA + a new Odia token + LID-CTC head + balanced sampling + non-speech data) for this specific seven-language set, a careful
speaker-disjoint evaluation against strong baselines, and the diagnosis and partial repair of the Indian-English regression.
\end{vivatip}
